Financial research code fails silently: a program that applies a stock split twice, or trades
on the close it has just read, still runs and still produces a backtest. We present \lib{}, an
agent-native Python library that gives a language-model agent what it needs to get such code
right: \emph{skills} (\NSkills{} dated skill documents, \NAlgorithms{} algorithms and
\NCards{} rule cards), \emph{checks} (\NGuards{} guards, among them \FixVisibleChecks{} property
checks that need no expected output) and \emph{self-improvement} (\NMethods{} methods under one
interface and one metered budget). It comes with \NTasks{} financial programming tasks, each
graded against an independent reference by a grader the agent never sees. On held-out tasks
that a single library call solves, \FixModels{} open-weight models were valid in
\FixFZeroPct\% of episodes without the library. Shown its interface, they mostly rewrote the
rule themselves (\FixFOnePct\%). When \texttt{route()} let them name the matching call first,
\FixFFourPct\% were valid, close to the \FixFTwoPct\% reached when we named it for them. Models
of \FixSizeCut{} parameters or more went from \FixLargeFZeroPct\% to \FixLargeFFourPct\%;
smaller ones never passed \FixSmallMaxPct\%. Without the library, the self-improvement
methods helped only models that already solved some tasks in a single attempt. These results are
exploratory,
since the library was revised on the same tasks before the final run. The library, the tasks
and every run record will be released.
